\documentclass{article}
\usepackage[utf8]{inputenc}
\usepackage[margin=1in,left=1.5in,includefoot]{geometry}
\usepackage{booktabs}
\usepackage{float}
% Header & Footer Stuff

\usepackage{fancyhdr}
\usepackage{wasysym}
\usepackage{amsmath}
\usepackage{tikz}
\usepackage{tabularx}
\usepackage{amssymb}
\usepackage{graphicx}
\pagestyle{fancy}
\lhead{MATH 220/199}
\rhead{March 10, 2026}
\fancyfoot{}
\lfoot{David Gallardo}
\fancyfoot[R]{}
%027601193
% The Main Document
\begin{document}

\begin{center}
 \LARGE\bfseries MATH 220/199 Week 8 Day 1
 \line(1,0){430}
\end{center}
\section*{Reminders}
\begin{enumerate}
    \item Office hours are held in Davenport 330 at 3pm on Thursdays.
    \item Do the Merit Quiz on the Merit Canvas page.
\end{enumerate}
\section*{We can do Infinitely Better than Linearization}
Linearization works because we have a reference point and understand the nearby behavior of a graph near the point of interest - What if we also knew about the behavior of the behavior of the graph? Let's estimate $(2.01)^3$
\begin{enumerate}
    \item First, using linearization, find an estimate for $(2.01)^3$. Call your linearization equation $L_1(x)$.
    \vfill
    \item Define $L_2(x) = L_1(x)+\frac{f''(a)}{2!}(x-a)^2 = f(a)+f'(a)f'(x-a)+\frac{f''(a)}{2!}(x-a)^2$. Use $L_2(x)$ to estimate $(2.01)^3$. Is this estimate better?
    \vfill
    \item Define $L_3(x) = L_2(x)+\frac{f'''(a)}{3!}(x-a)^3$ and re-estimate $(2.01)^3$
    \vfill
    \item In general, we can construct these equations which give us better and better estimates using the recursive formula $L_n(x) = L_{n-1}(x)+\frac{f^{(n)}(a)}{n!}(x-a)^n$ - in Calc 2 you will call this a "Taylor Series". If you have time, compute $L_4(2.01)$ and see how good this estimate is.
    \vfill
\end{enumerate}
\newpage
\section*{Plausible Exam Word Problem}
Suppose that a particle oscillates and has position function given by $p(t) = \sin(t)-\cos(t)$
\begin{enumerate}
    \item Over the interval $[0, 3\pi]$, determine all the times $t$ when the particle is not moving.
    \vfill
    \item Over the interval $[0, \pi]$, determine the intervals where the particle is speeding up, and where it is slowing down.
    \vfill
    \item Over the interval $[0, \pi]$, determine when the particle has the greatest speed.
    \vfill
    \item The change in acceleration over time is called 
    "jerk". If $v(t)$ is the velocity equation, find $j(t)$, the jerk equation in terms of $v(t)$.
    \vfill
\end{enumerate}

\end{document}
